\documentclass[letterpaper]{article}
\usepackage{amsmath,amssymb}
\usepackage[letterpaper,margin=1in]{geometry}
\setlength{\parindent}{0pt}
\newcommand{\noteheader}[1]{\par\bigskip\noindent\textbf{#1}\par\medskip}

\begin{document}

\textbf{CS 2050 --- Notes 09/23}

\noteheader{Sets and elements}

A \textbf{set} is an unordered collection of distinct objects. The objects
in a set are its \textbf{elements}. For a set $A$ and an object $a$,
\[
\begin{array}{ll}
a\in A & a\text{ is an element of }A,\\[2pt]
a\notin A & a\text{ is not an element of }A.
\end{array}
\]
Uppercase letters usually name sets; lowercase letters usually name elements.
Two sets are equal if and only if they have the same elements.

\noteheader{Set-builder notation}

\[
\{x\mid x\text{ has property }P\}
\]
denotes the set of all $x$ with property $P$.

\noteheader{Common number sets}

\[
\begin{array}{ll}
\mathbb N=\{0,1,2,\ldots\} & \text{natural numbers},\\[3pt]
\mathbb Z=\{\ldots,-2,-1,0,1,2,\ldots\} & \text{integers},\\[3pt]
\mathbb Z^+=\{1,2,\ldots\} & \text{positive integers},\\[3pt]
\mathbb Q=\{p/q\mid p,q\in\mathbb Z,\ q\ne0\} & \text{rational numbers},\\[3pt]
\mathbb R^+ & \text{positive real numbers},\\[3pt]
\mathbb C & \text{complex numbers}.
\end{array}
\]

\clearpage

\noteheader{Empty and singleton sets}

The \textbf{empty set} has no elements and is denoted by $\varnothing$ or
$\{\}$. A \textbf{singleton set} has exactly one element.

\noteheader{Universal set}

In a Venn diagram, the \textbf{universal set} $U$ contains all objects under
consideration.

\noteheader{Subsets and supersets}

$A$ is a \textbf{subset} of $B$, written $A\subseteq B$, if every element
of $A$ is also an element of $B$. Equivalently, $B$ is a \textbf{superset}
of $A$, written $B\supseteq A$.
\[
A\subseteq B\quad\Longleftrightarrow\quad
\forall x\,(x\in A\to x\in B).
\]
To show $A\not\subseteq B$, find an element of $A$ that is not in $B$:
\[
\exists x\,(x\in A\land x\notin B).
\]

\clearpage

\noteheader{Theorem}

For every set $S$, $\varnothing\subseteq S$ and $S\subseteq S$.

\noteheader{Proof}

To show $\varnothing\subseteq S$, consider the statement
$x\in\varnothing\to x\in S$. There is no $x\in\varnothing$, so the
implication is vacuously true for every $x$. Also, $x\in S\to x\in S$
is true for every $x$, giving $S\subseteq S$.

\noteheader{Proper subsets}

$A$ is a \textbf{proper subset} of $B$, written $A\subsetneq B$, when
$A\subseteq B$ and $A\ne B$.

\noteheader{Cardinality}

If a set $S$ contains exactly $n$ distinct elements, where $n$ is a
nonnegative integer, then $S$ is \textbf{finite} and its
\textbf{cardinality} is $|S|=n$. For example,
\[
|\mathbb Z|\text{ is infinite},
\qquad
\bigl|\{\varnothing,\mathbb N,\mathbb Z\}\bigr|=3.
\]

\noteheader{Power sets}

The \textbf{power set} of $S$ is the set of all subsets of $S$, denoted
$\mathcal P(S)$. If $|S|=n$, then
\[
|\mathcal P(S)|=2^n.
\]

\clearpage

\noteheader{Cartesian product}

For sets $A$ and $B$, their \textbf{Cartesian product} is the set of all
ordered pairs whose first component is in $A$ and second component is in $B$:
\[
A\times B=\{(a,b)\mid a\in A\land b\in B\}.
\]

\noteheader{Example}

The product of any set $A$ with the empty set is empty:
\begin{align*}
A\times\varnothing
&=\{(a,b)\mid a\in A\land b\in\varnothing\}\\
&=\varnothing,
\end{align*}
because no choice of $b\in\varnothing$ is possible.

\end{document}
